\documentclass[12pt]{article}

% ── Packages ──────────────────────────────────────────────────────────────────
\usepackage[margin=1in]{geometry}
\usepackage{booktabs}
\usepackage{graphicx}
\usepackage{hyperref}
\usepackage{natbib}
\usepackage{setspace}
\usepackage{amsmath}
\usepackage{microtype}

\onehalfspacing
\bibpunct{(}{)}{;}{a}{,}{,}

% ── Title ─────────────────────────────────────────────────────────────────────
\title{% TODO: write your title here
  Replication and Extension of Magistro et al. (2026):\\
  % TODO: add a subtitle that captures your extension
  {[Your Extension Title Here]}
}

\author{
  % TODO: add your names
  Author One \and Author Two
}

\date{\today}

% ── Document ──────────────────────────────────────────────────────────────────
\begin{document}

\maketitle

% ── Abstract ──────────────────────────────────────────────────────────────────
\begin{abstract}
% TODO: 3–4 sentences.
% Sentence 1: What question does this paper address?
% Sentence 2: What is the design and data?
% Sentence 3: What do you find?
% Sentence 4 (optional): What does it imply?
[Abstract placeholder.]
\end{abstract}

\newpage

% ── 1. Introduction ───────────────────────────────────────────────────────────
\section{Introduction}
\label{sec:intro}

% TODO: ~3–4 paragraphs.
% Para 1: Motivate the question. Why do public attitudes toward AI-driven
%         economic disruption matter for politics?
% Para 2: What does Magistro et al. 2026 find? Briefly summarize the
%         paper's claim, design, and main result.
% Para 3: What does your extension add? State the extension question clearly
%         and explain why it matters given the paper's argument.
% Para 4 (optional): Road map sentence.

[Introduction placeholder.]

% ── 2. Data and Design ────────────────────────────────────────────────────────
\section{Data and Design}
\label{sec:design}

% TODO: ~2 paragraphs.
% Para 1: Describe the conjoint experiment — sample (N=6,056, US+Canada, YouGov),
%         conjoint attributes (price, customer service jobs, factory jobs,
%         data science jobs), treatment arms (AI vs. Offshoring framing),
%         and outcome (support rating 1–5).
% Para 2: Describe your extension approach — what additional variable or
%         specification you introduce and why it addresses your question.

[Data and design placeholder.]

% ── 3. Results ────────────────────────────────────────────────────────────────
\section{Results}
\label{sec:results}

% TODO: ~2–3 paragraphs + figure/table.
% Para 1: Replication — summarize what Figure~\ref{fig:replication} shows.
%         Does your Figure 1 match Magistro et al.?
% Para 2: Extension — describe your main finding.
%         Reference your figure or table (Figure~\ref{fig:extension} or
%         Table~\ref{tab:main}).
% Para 3 (optional): Any surprising or null result worth noting.

\subsection{Replication}

[Replication results placeholder.]

\begin{figure}[h]
  \centering
  \includegraphics[width=0.85\textwidth]{figures/figure1_replication.png}
  \caption{Replication of Figure 1: Marginal means by treatment condition
           (AI vs. Offshoring). Points show marginal means; horizontal lines
           are 95\% confidence intervals. Dashed vertical line at 3 (midpoint
           of the 1--5 support scale).}
  \label{fig:replication}
\end{figure}

\subsection{Extension}

[Extension results placeholder.]

\begin{figure}[h]
  \centering
  \includegraphics[width=0.85\textwidth]{figures/extension_result.png}
  \caption{% TODO: write a caption that describes what this figure shows.
  }
  \label{fig:extension}
\end{figure}

% Uncomment if you have a regression table:
% \input{tables/main.tex}

% ── 4. Discussion ─────────────────────────────────────────────────────────────
\section{Discussion}
\label{sec:discussion}

% TODO: ~2 paragraphs.
% Para 1: What do your results mean for the paper's argument about common
%         microfoundations of AI and offshoring attitudes?
% Para 2: Limitations and what a future study would need.

[Discussion placeholder.]

% ── 5. Conclusion ─────────────────────────────────────────────────────────────
\section{Conclusion}
\label{sec:conclusion}

% TODO: 2–3 sentences. Restate the contribution; no new information.

[Conclusion placeholder.]

% ── References ────────────────────────────────────────────────────────────────
\bibliographystyle{apsr}
\bibliography{refs}

\end{document}
